\title{FlashAttention-3: Fast and Accurate Attention with Asynchrony and Low-precision}

\begin{abstract}
As the principal component of the widespread Transformer framework, attention serves as the limiting factor for scaling large language models and processing extended sequences. \textsc{FlashAttention} previously introduced a methodology to accelerate attention mechanisms on GPUs by reducing memory access costs. Nonetheless, this approach does not yet fully exploit recent hardware advancements; for instance, \textsc{FlashAttention-2} manages only 35\% hardware utilization on the H100 GPU. In this work, we propose three strategies to further improve attention speed on Hopper GPUs: leveraging the concurrent operation of Tensor Cores and TMA to (1) mask memory latency and computation through warp specialization, (2) alternate between block-level matrix multiplication and softmax operations, and (3) employ block-based quantization along with noncoherent workflows that utilize FP8 low-precision support provided by hardware. Our empirical analysis reveals that \textsc{FlashAttention-3} achieves a performance gain of 1.5-2.0$\times$ on H100 GPUs—reaching up to 740 TFLOPs/s in FP16 mode (which corresponds to 75\% utilization) and nearly 1.2 PFLOPs/s in FP8. Moreover, we confirm that FP8 \textsc{FlashAttention-3} yields 2.6$\times$ less numerical error compared to a naïve FP8 attention approach.
\end{abstract}

\section{Introduction}
\label{sec:intro}

Within the Transformer framework~\citep{vaswani2017attention}, the primary limitation on computational efficiency arises from the attention mechanism, as calculating self-attention scores for queries and keys exhibits a quadratic growth with respect to sequence length. Enabling attention operations over extended contexts could facilitate a range of advanced capabilities—including improved modeling and reasoning across multiple lengthy documents~\citep{guo2021longt5,shaham2022scrolls,peng2023yarn} or extensive code repositories~\citep{roziere2023code, li2023starcoder}—as well as supporting new data types, such as high-resolution images~\citep{chen2022scaling}, audio~\citep{gulati2020conformer}, and video~\citep{ho2022video}. Moreover, such advancements enable novel applications in domains like recommendation with lengthy interaction histories~\citep{sun2019bert4rec} and agents tasked with long-horizon workflows~\citep{yao2022react}. This has sparked a burgeoning research effort aimed at accelerating attention for longer input sequences, leveraging approximate computation~\citep{katharopoulos2020transformers,choromanski2020rethinking,
tay2020efficient}, optimizations at the software level (\citep{rabe2021self,dao2022flashattention,kwon2023efficient}), and even the design of fundamentally new model architectures~\citep{peng2023rwkv,sun2023retentive,gu2023mamba}.

In this study, we extend the research by \citet{dao2022flashattention}, who designed exact-attention algorithms informed by the GPU’s computational paradigm and architectural details at the conceptual level. Dao et al., in \citep{dao2022flashattention}, proposed \textsc{FlashAttention}, introducing an innovative tiling approach that enables efficient parallelization of attention by combining all operations into a unified GPU kernel and thereby eliminating the need for intermediate global memory accesses. Subsequently, \citet{dao2023flashattention2} presented a redesigned version named \textsc{FlashAttention-2}, which further enables parallelization across the sequence dimension and restructures the forward pass’s innermost computation to operate on blocks of key and value matrices. This leads to better load balancing and higher GPU occupancy. Despite these advancements, our analysis indicates that \textsc{FlashAttention-2} remains suboptimal with respect to hardware utilization on the latest generation GPUs; for example, it achieves utilization rates of only around 35\% compared to the 80-90\% observed for highly optimized GEMM kernels on the Hopper H100 GPU. This discrepancy is due in part to details of the implementation—such as the failure to replace Ampere-centric instructions with Hopper-specific alternatives for Tensor Core usage. Recent efforts such as ThunkerKitten~\citep{spector2024thunder} and the release of cuDNN 9~\citep{cudnn9} have demonstrated that leveraging Hopper’s specialized instructions and employing tile-based methods can yield significant acceleration of attention mechanisms while also reducing implementation complexity.

At a foundational level, the algorithm underlying \textsc{FlashAttention-2} operates according to a straightforward, synchronous paradigm, lacking integration of asynchrony and low-precision as core design principles. Asynchrony arises due to the division of labor in specialized hardware components optimized for key ML tasks—for instance, matrix operations are accelerated on dedicated units like Tensor Cores, while data transfers are handled by components such as the Tensor Memory Accelerator (TMA), both functioning independently of the conventional CUDA cores that manage general computation, logic, and arithmetic operations. The adoption of low-precision formats, including FP8 in Hopper and FP4 in Blackwell—following the adoption of FP16 (Pascal, 2017) and BF16 (Ampere, 2020)—has demonstrated marked improvements in computational throughput, delivering two to four times more output per unit power and silicon area. \cref{subsec:hardware} provides an overview of how Hopper supports these advanced hardware features. The principal technical obstacle involves reengineering \textsc{FlashAttention-2} to leverage these advantages: exploiting asynchrony means orchestrating matmul and softmax to execute concurrently, despite data dependencies, while effective utilization of low-precision demands strategies for minimizing quantization artifacts, particularly concerning the handling of rare or atypical features present in LLMs~\citep{dettmers2208llm, sun2024massive}.

Accordingly, we introduce \textsc{FlashAttention-3}, which integrates three novel techniques designed to enhance efficiency on advanced GPU platforms:\footnote{While our findings are presented with reference to NVIDIA's Hopper architecture, the algorithm is applicable to any GPU system that supports strong asynchronous operations and efficient low-precision computation.}

\begin{enumerate}

\item \textbf{Producer-Consumer asynchrony:} We introduce a warp-specialized variant of software pipelining in which data producers and consumers are mapped to distinct warps, capitalizing on the decoupled operation of data transfers and computations on Tensor Cores. This design extends the algorithm’s capacity to mask both memory latency and instruction dispatch delays.
\item \textbf{Hiding softmax under asynchronous block-wise GEMMs:} The inherently lower-throughput softmax-related operations—such as floating point multiply-add and exponentials—are overlapped with asynchronous execution of block GEMMs using WGMMA instructions. To realize this, we modify the \textsc{FlashAttention-2} algorithm, eliminating specific sequential dependencies between the softmax computation and GEMMs. For instance, in the two-stage scheme, while the softmax acts on one block of the score matrix, asynchronous WGMMA prefetches and processes the next block.
\item \textbf{Hardware-accelerated low-precision GEMM:} We modify the forward pass to enable the use of FP8 Tensor Cores for GEMM operations, yielding nearly a 2$\times$ gain in TFLOPs/s throughput. This adjustment involves reconciling differences in WGMMA’s memory layout expectations for FP8 operands and FP32 accumulators. Techniques such as block quantization and incoherent processing are utilized to alleviate accuracy degradation due to reduced precision with FP8.
\end{enumerate}

For empirical validation, we conduct benchmarking of \textsc{FlashAttention-3} on the H100 SXM5 GPU across diverse settings. The results demonstrate that (1) FP16 realizes a 1.5-2.0$\times$ acceleration over \textsc{FlashAttention-2} during the forward pass (with throughput reaching 740 TFLOPs/s) and achieves a 1.5-1.75$\times$ speed increase in the backward pass; (2) FP8 attains nearly 1.2 PFLOPs/s; and (3) for large sequence lengths, FP16 delivers superior performance, while FP8 shows competitive results\footnote{Specifically, for a head dimension of 64, \textsc{FlashAttention-3} FP8 surpasses other implementations, whereas for head dimensions 128 and 256 it is on par without causal masking, and lags behind when causal masking is applied.} compared to the state-of-the-art attention kernel provided by NVIDIA’s cuDNN library. Additionally, we establish that the FP16 version of \textsc{FlashAttention-3} maintains the same numerical error profile as \textsc{FlashAttention-2} and surpasses the conventional attention mechanism, attributed to the retention of intermediate results (such as softmax rescaling) in FP32. Furthermore, FP8 \textsc{FlashAttention-3} employing block quantization and incoherent processing achieves 2.6$\times$ higher accuracy than the baseline attention approach using per-tensor quantization in scenarios where outlier features are present.

We release \textsc{FlashAttention-3} under a permissive license\footnote{\textsc{FlashAttention-3}
  is available at \texttt{https://github.com/Dao-AILab/flash-attention}}, and intend to incorporate it into the PyTorch and Hugging Face ecosystems to maximize its accessibility and impact within the research and developer communities.

\section{Background: Multi-Head Attention and GPU Characteristics}
\label{sec:background}

\subsection{Multi-Head Attention}
\label{subsec:multi_head_attn}

Let $\mathbf{Q}, \mathbf{K}, \mathbf{V} \in \mathbb{R}^{N \times d}$ be the query, key and value input sequences associated to a single head, where $N$ is the sequence length and $d$ is the head dimension. Then the attention output $\mathbf{O}$ is computed as:

\begin{equation*}
  \mathbf{S} = \alpha \mathbf{Q} \mathbf{K}^\top \in \mathbb{R}^{N \times N}, \quad \mathbf{P} = \mathrm{softmax}(\mathbf{S}) \in \mathbb{R}^{N \times N}, \quad \mathbf{O} = \mathbf{P}\mathbf{V} \in \mathbb{R}^{N \times d},
\end{equation*}

In this formulation, the $\mathrm{softmax}$ function is executed on a per-row basis, and the scaling parameter $\alpha$ is commonly chosen as $1/\sqrt{d}$. To enhance numerical stability during exponentiation, it is standard practice to subtract $\mathrm{rowmax}(\mathbf{S})$ from $\mathbf{S}$. Within the multi-head attention (MHA) mechanism, each head is assigned distinct projections for queries, keys, and values. These operations are conducted in parallel for all heads and over batches, resulting in the complete output tensor.

Now let $\phi$ be a scalar loss function and let $\mathbf{d}(-) = \partial \phi / \partial (-)$ be notation for the gradient. Given the output gradient $\mathbf{dO} \in \mathbb{R}^{N \times d}$, we compute $\mathbf{dQ}$, $\mathbf{dK}$, and $\mathbf{dV}$ according to the chain rule as follows:

\begin{align*}
  \mathbf{dV} &= \mathbf{P}^\top \mathbf{dO} \in \mathbb{R}^{N \times d} \\
  \mathbf{dP} &= \mathbf{dO} \mathbf{V}^\top \in \mathbb{R}^{N \times N} \\
  \mathbf{dS} &= \mathrm{dsoftmax} (\mathbf{dP}) \in \mathbb{R}^{N \times N} \\
  \mathbf{dQ} &= \alpha \mathbf{dS} \mathbf{K} \in \mathbb{R}^{N \times d} \\
  \mathbf{dK} &= \alpha \mathbf{dS}^\top \mathbf{Q} \in \mathbb{R}^{N \times d},
\end{align*}

In this context, we obtain $\mathbf{d}s = (\mathrm{diag}(p) - p p^\top)\mathbf{d}p$ where $p = \mathrm{softmax}(s)$, with $s$ being a vector, and denote the row-wise application of this expression as $\mathrm{dsoftmax}(\mathbf{dP})$. This operation, in turn, supports efficient parallel computation over both heads and batch elements when executing the backward pass in MHA.

\subsection{GPU hardware characteristics and execution model}
\label{subsec:hardware}

We outline the facets of the GPU execution model pertinent to \textsc{FlashAttention-3}, concentrating specifically on how these are realized within the NVIDIA Hopper architecture.

\paragraph{Memory hierarchy:} The memory architecture of the GPU consists of several levels, where smaller, faster memories provide higher bandwidth, as shown in \cref{tab:gpu-hierarchy}\footnote{\citet{luo2024benchmarking} gives a shared memory bandwidth of 128 bytes per clock cycle per SM; multiplying this by 132 SMs and a boost clock of 1830 MHz yields the reported aggregate bandwidth.}. Global memory (GMEM)—often referred to as HBM—serves as the off-chip DRAM and is available to all streaming multiprocessors (SMs). Data originating from GMEM is automatically stored in an on-chip L2 cache before further processing. Each SM is also equipped with a limited-size, programmer-controlled, highly banked shared memory (SMEM) located on-chip. At the closest level to computation, each SM has its own register file.

\paragraph{Thread hierarchy:} Within the GPU programming paradigm, execution units are systematically organized into a hierarchy of threads. This structure begins with individual threads at the lowest level, then aggregates these into warps, which each contain 32 threads. Next, four adjacent warps form a warpgoup, followed by threadblocks (also known as cooperative thread arrays or CTAs). At a higher level, the Hopper architecture introduces threadblock clusters, with the largest grouping in this hierarchy referred to as grids.

There is a strong relationship between these two hierarchical structures. Threads belonging to a single CTA are assigned simultaneously to a single SM, while multiple CTAs that are part of the same cluster are concurrently dispatched to the same GPC. The SMEM memory space can be accessed by all threads of a CTA, while each individual thread is limited to a maximum of 256 private registers (RMEM).

\paragraph{Asynchrony and warp-specialization:}

GPUs function as throughput-oriented processors that depend on both parallelism and asynchronous operations to mask delays associated with memory access and instruction execution. To support asynchronous memory transfers between GMEM and SMEM, Hopper introduces a specialized hardware component known as the Tensor Memory Accelerator (TMA) \cite[\S7.29]{cuda}. Additionally, distinguishing itself from earlier architectures like Ampere, Hopper's Tensor Core is accessible through the warpgroup-wide WGMMA instruction \cite[\S9.7.14]{ptx} and is designed for asynchronous operation, with the capability to directly retrieve inputs from shared memory.

Asynchronous hardware capabilities enable the use of warp-specialized kernels, wherein the warps within a CTA are designated exclusively as either producers (handling data movement) or consumers (performing computation), with each role restricted to its respective operation. This separation generally enhances the compiler’s effectiveness in producing efficient instruction schedules \citep{warp-specialization-2011}. Moreover, Hopper architecture facilitates the dynamic redistribution of registers among warpgroups using the \verb|setmaxnreg| command \cite[\S9.7.17.1]{ptx}, allowing warps assigned to MMAs to utilize a greater portion of RMEM compared to those restricted to TMA operations, which typically require only one active thread.

\paragraph{Low-precision number formats:}
\label{sec:low-precision-gpu}
Current GPUs are equipped with dedicated hardware to boost the performance of low-precision arithmetic operations. As an illustration, Hopper’s FP8 Tensor Cores can be utilized via the WGMMA instruction, enabling each SM to achieve double the throughput relative to FP16 or BF16 operations.

Proper utilization of FP8 WGMMA requires familiarity with the operand layout limitations. Consider a GEMM operation computing $A \times B^{\top}$ where $A$ is of size $M\times K$ and $B$ is $N\times K$. We define an operand as \emph{mn-major} if its elements are contiguous along the outer dimension ($M$ for $A$, $N$ for $B$), and as \emph{k-major} if continuity is along the inner $K$ dimension. With FP16 WGMMA, either mn-major or k-major layouts in SMEM are permissible for operands; in contrast, FP8 WGMMA accepts only the k-major configuration. Furthermore, in use cases such as attention—where the goal is to combine sequential GEMM operations within one kernel—mismatches between the memory layouts of FP32 accumulators and FP8 operands can hinder the execution of subsequent dependent FP8 WGMMAs.

Regarding attention mechanisms, these layout limitations require specific adaptations to be made to the FP8 algorithm’s design, as discussed in detail in \cref{sec:algofp8}.

\subsection{Standard Attention and Flash Attention} As described by \citet{dao2022flashattention}, \textbf{standard attention} refers to the conventional GPU-based approach in which the intermediate matrices $\mathbf{S}$ and $\mathbf{P}$ are explicitly stored in HBM. \textsc{FlashAttention} was introduced to optimize this process by utilizing a localized softmax reduction, thereby eliminating costly intermediate memory operations and merging the attention computation into a single kernel. The localized softmax is realized in lines \ref{code-ws:softmax_start}-\ref{code-ws:softmax_end} of the consumer mainloop within \cref{alg:flash3_wgmma_ws_only}, along with appropriate scaling of $\mathbf{O}$ blocks. A straightforward derivation verifying that this procedure accurately produces $\mathbf{O}$ is presented in \cite[\S 2.3.1]{dao2023flashattention2}.

\section{FlashAttention-3: Algorithm}
\label{sec:algo}

This section presents the \textsc{FlashAttention-3} algorithm. Our discussion centers on the forward pass for clarity, while details regarding the backward pass are deferred to~\cref{sec:algo_ws_bwd}. We begin by demonstrating the incorporation of warp-specialization and the use of a circular SMEM buffer within the foundational \textsc{FlashAttention-2} framework. Next, we detail the utilization of WGMMA’s asynchronous properties to establish a 2-stage pipeline that overlaps the GEMM and softmax computations. We conclude with a discussion of the adaptations required to support FP8, including adjustments for layout compatibility and enhancements to accuracy through block quantization and incoherent computation techniques.

\subsection{Producer-Consumer asynchrony through warp-specialization and pingpong scheduling}
\label{sec:algo_ws}

\paragraph{Warp-specialization}
Similar to \textsc{FlashAttention-2}, the forward propagation in \textsc{FlashAttention-3} admits substantial parallelism across the batch dimension, the set of attention heads, and the length of the query sequence. Consequently, presenting the algorithm from a CTA-level perspective is sufficient, where each CTA processes a specific tile $\mathbf{Q}_i$ from the query matrix to generate the corresponding segment $\mathbf{O}_i$ of the output. For clarity, the warp-specialization strategy is initially described using a circular SMEM buffer, intentionally omitting any overlapping between GEMM and softmax stages. Denote the head dimension as $d$, sequence length as $N$, and select a block size $B_r$ for the queries so that $\mathbf{Q}$ can be partitioned into $T_r = \lceil \frac{N}{B_r} \rceil$ blocks, specifically $\mathbf{Q}_1, .., \mathbf{Q}_{T_r}$.

\begin{algorithm}[H]
    \caption{\small\label{alg:flash3_wgmma_ws_only}\textsc{FlashAttention-3} forward pass \textbf{excluding} intra-consumer overlapping -- CTA perspective}
    \begin{algorithmic}[1]
\REQUIRE Input matrices $\mathbf{Q}_i \in \mathbb{R}^{B_r \times d}$, $\mathbf{K}, \mathbf{V} \in \mathbb{R}^{N \times d}$ residing in HBM; key block size $B_c$; number of key blocks $T_c = \lceil \frac{N}{B_c} \rceil$.
\STATE Create and set up a pipeline structure that coordinates barrier synchronization across an $s$-stage circular SMEM buffer.
\IF {operating as a producer warpgroup}
\STATE Free up a set number of registers as predetermined.
\STATE Execute data load of $\mathbf{Q}_i$ from HBM into shared memory.
\STATE After load completes, trigger commit to signal consumers regarding completion of $\mathbf{Q}_i$ load.
\FOR{$0 \le j < T_c$}
    \STATE Wait until the buffer's $(j\,\%\,s)$th stage becomes available from consumption.
    \STATE Fetch blocks $\mathbf{K}_j$ and $\mathbf{V}_j$ from HBM into the corresponding buffer stage in shared memory.
    \STATE Following load completion, commit to inform consumers of $\mathbf{K}_j$ and $\mathbf{V}_j$ availability.
\ENDFOR
\ELSE \STATE Allocate back the specified number of registers in proportion to the quantity of consumer warps.
\STATE Within chip memory, initialize $\mathbf{O}_i = (0) \in \mathbb{R}^{B_r \times d}$, alongside $\ell_i, m_i = (0), (-\infty) \in \mathbb{R}^{B_r}$.
\STATE Stall until $\mathbf{Q}_i$ becomes accessible in shared memory.
\FOR{$0 \le j < T_c$}
\STATE Wait for the current key block $\mathbf{K}_j$ to be available in shared memory.
\STATE Carry out computation $\mathbf{S}_i^{(j)} = \mathbf{Q}_i \mathbf{K}_j^T$ (using SS-GEMM); commit afterward and wait. \label{alg:ws_only_gemm1}
\STATE Save $m_i^{\mathrm{old}} = m_i$; update $m_i = \mathrm{max}(m_i^{\mathrm{old}}, \mathrm{rowmax}(\mathbf{S}_i^{(j)}))$. \label{code-ws:softmax_start}
\STATE Compute $\widetilde{\mathbf{P}}_i^{(j)} = \mathrm{exp}(\mathbf{S}_i^{(j)} - m_i)$ and adjust $\ell_i = \mathrm{exp}(m_i^{\mathrm{old}} - m_i) \ell_i + \mathrm{rowsum}(\widetilde{\mathbf{P}}_i^{(j)})$. \label{code-ws:softmax_end}
\STATE Wait for shared memory load of $\mathbf{V}_j$ to finish.
\STATE Perform update $\mathbf{O}_i = \mathrm{diag}(\mathrm{exp}(m_i^{\mathrm{old}} - m_i))^{-1} \mathbf{O}_i + \widetilde{\mathbf{P}}_i^{(j)} \mathbf{V}_j$ (RS-GEMM); commit changes and await next step. \label{alg:ws_only_gemm2}
\STATE Release ownership of the $(j\,\%\,s)$th buffer stage back to the producer.
\ENDFOR
\STATE Apply normalization: $\mathbf{O}_i = \mathrm{diag}(\ell_i)^{-1} \mathbf{O}_i$ and set $L_i = m_i + \log(\ell_i)$.
\STATE Output results by storing $\mathbf{O}_i$ and $L_i$ in HBM as the $i$th segment of $\mathbf{O}$ and $L$.
\ENDIF
\end{algorithmic}
\end{algorithm}

In our deployment of \cref{alg:flash3_wgmma_ws_only} on the Hopper architecture, we make use of \verb|setmaxnreg| for the purpose of (de)allocating resources, TMA to handle loading of $\mathbf{Q}_i$ as well as the sets $\{ \mathbf{K}_j, \mathbf{V}_j \}_{0 \leq j < T_c}$, and WGMMA to perform GEMM operations within the consumer mainloop. Here, the SS and RS prefixes specify whether the primary input operand originates from shared memory or the register file, respectively. When analyzing the execution sequence in \cref{alg:flash3_wgmma_ws_only}, it is important to observe that TMA load instructions function asynchronously and thus are not hindered by previous outstanding loads. Additionally, during the initial $s$ iterations of the producer mainloop—while the buffer is being populated—no waits are executed.

\paragraph{Pingpong scheduling} The asynchronous execution model afforded by WGMMA and TMA, coupled with warp-specialization strategies, enables overlapping the softmax phase carried out by one warpgroup with the concurrent GEMM operations of a different warpgroup. This design is motivated by the observation that on modern hardware accelerators, the peak throughput for matrix multiplication operations substantially exceeds that of non-matmul tasks. For example, on the H100 SXM5 GPU, the maximum theoretical throughput for FP16 matrix multiplication operations reaches 989 TFLOPS, while special function computations, such as exponentials required in softmax, reach only 3.9 TFLOPS\footnote{According to the CUDA programming guide, each streaming multiprocessor (SM) can execute 16 special function operations per cycle. Multiplying 16 by 132 SMs and a clock speed of 1830 MHz results in 3.9 TFLOPS for these functions.}. In the FP16 attention forward pass with a head dimension of 128, the number of FLOPs for matrix multiplications exceeds those for exponentials by a factor of 512. However, the hardware delivers 256 times less throughput for exponentials, resulting in softmax taking up to half the total compute cycles as matrix multiplication. In FP8 precision, this disparity intensifies because the matrix multiplication throughput is doubled, whereas the rate for exponentials remains constant.

Because the exponential computation is handled by a dedicated hardware component (the multi-function unit), it is optimal to schedule this calculation concurrently with matmul operations on the Tensor Cores. To achieve this coordination, we employ synchronization barriers (\texttt{bar.sync} instructions), which ensure that the GEMMs (specifically, GEMM1 -- $\mathbf{P} \mathbf{V}$ of the current iteration and GEMM0 -- $\mathbf{Q} \mathbf{K}^\top$ from the next iteration) of warpgroup 1 execute before the GEMMs in warpgroup 2. Consequently, the softmax computation for warpgroup 1 is carried out simultaneously as warpgroup 2 is executing its GEMMs. The process then alternates: warpgroup 2 proceeds with the softmax operation while warpgroup 1 is engaged in GEMMs, a scheme we refer to as ``pingpong'' scheduling. This mechanism is depicted in~\cref{fig:pingpong_scheduling}. 

Although, in real-world scenarios, the pingpong scheduling does not perfectly match the conceptual depiction in the figure, our observations indicate that this approach typically yields better performance. For example, we observe improvements from 570 TFLOPS to between 620 and 640 TFLOPS for FP16 forward passes with a head dimension of 128 and sequence length of 8192.

\paragraph{Attention variants} For both multi-query attention~\citep{shazeer2019fast} and grouped query attention~\citep{ainslie2023gqa}, our implementation adheres to the method used in \textsc{FlashAttention-2}, modifying the tensor indexing scheme to prevent unnecessary duplication of $\mathbf{K}$ and $\mathbf{V}$ within HBM.

\subsection{Intra-warpgroup overlapping GEMMs and softmax}

It is possible to interleave certain instructions from the softmax computation with instructions from the GEMMs within a single warpgroup. Here, we outline one approach that enables such overlap.

In the attention algorithm, operations within the inner loop (main loop) have sequential dependencies that impede parallelization within a single iteration. For example, (local) softmax (lines \ref{code-ws:softmax_start} to \ref{code-ws:softmax_end}) relies on the output $\mathbf{S}_i^{(j)}$ of the first GEMM, while the second GEMM takes its result $\widetilde{\mathbf{P}}_i^{(j)}$ as an operand. Indeed, the wait statements in lines \ref{alg:ws_only_gemm1} and \ref{alg:ws_only_gemm2} of \cref{alg:flash3_wgmma_ws_only} serialize the execution of softmax and GEMMs. However, we can break these dependencies by pipelining across iterations through additional buffers in registers. Pursuing this idea, we propose the following two-stage\footnote{Note that the number of stages of the overlapping scheme is bounded by, but need not equal, the number $s$ of stages in the circular SMEM buffer.} GEMM-softmax pipelining algorithm:

\begin{algorithm}[H]
    \caption{\small\label{alg:flash3_wgmma}\textsc{FlashAttention-3} consumer warpgroup forward pass}
    \begin{algorithmic}[1]
      \REQUIRE  Matrices $\mathbf{Q}_i \in \mathbb{R}^{B_r \times d}$ and $\mathbf{K}, \mathbf{V} \in \mathbb{R}^{N \times d}$ located in HBM, key block dimension $B_c$, with $T_c = \lceil \frac{N}{B_c} \rceil$.
      \STATE Allocate the necessary number of registers in accordance with the number of consumer warps.
      \STATE Initialize on-chip values: set $\mathbf{O}_i = (0) \in \mathbb{R}^{B_r \times d}$, as well as $\ell_i = (0)$ and $m_i = (-\infty)$ within $\mathbb{R}^{B_r}$.
      \STATE Wait until both $\mathbf{Q}_i$ and $\mathbf{K}_0$ have been staged into shared memory.
      \STATE Using WGMMA, evaluate $\mathbf{S}_{\mathrm{cur}} = \mathbf{Q}_i \mathbf{K}_0^T$. Ensure commitment and synchronization.
      \STATE Free up the initial stage (0-th) of the $\mathbf{K}$ buffer.
      \STATE Update $m_{i}$, $\tilde{\mathbf{P}}_{\mathrm{cur}}$, and $\ell_{i}$ with respect to $\mathbf{S}_{\mathrm{cur}}$, and appropriately adjust the scale of $\mathbf{O}_i$.
      \FOR{$1 \le j < T_c - 1$}
        \STATE Wait for block $\mathbf{K}_j$ to be fetched into shared memory.
        \label{code:mainloop_start}
        \STATE Leverage WGMMA to compute $\mathbf{S}_{\mathrm{next}} = \mathbf{Q}_i \mathbf{K}_{j}^T$. Commit computation, but defer waiting. \label{code:first_wgmma}
        \STATE Wait until $\mathbf{V}_{j-1}$ is present in shared memory.
        \STATE Through WGMMA, calculate $\mathbf{O}_{i} = \mathbf{O}_{i} + \tilde{\mathbf{P}}_{\mathrm{cur}} \mathbf{V}_{j-1}$ and commit immediately, bypassing any wait. \label{code:second_wgmma}
        \STATE Wait for the WGMMA operation $\mathbf{Q}_i \mathbf{K}_{j}^T$ to complete.
        \STATE Derive $m_{i}$, $\tilde{\mathbf{P}}_{\mathrm{next}}$, and $\ell_{i}$ based on the newly computed $\mathbf{S}_{\mathrm{next}}$. \label{code:softmax}
        \STATE Once WGMMA on $\tilde{\mathbf{P}}_{\mathrm{cur}} \mathbf{V}_{j-1}$ completes, perform rescaling of $\mathbf{O}_i$.
        \STATE Release the $(j\,\%\,s)$-th stage for $\mathbf{K}$ and the $(j-1\,\%\,s)$-th stage for $\mathbf{V}$ in their respective buffers.
        \STATE Copy $\mathbf{S}_{\mathrm{next}}$ contents over to $\mathbf{S}_{\mathrm{cur}}$.
        \label{code:mainloop_end}
      \ENDFOR
      \STATE Await completion of loading for $\mathbf{V}_{T_c - 1}$ into shared memory.
      \STATE Use WGMMA to calculate 
      $\mathbf{O}_{i} = \mathbf{O}_{i} + \tilde{\mathbf{P}}_{\mathrm{last}} \mathbf{V}_{T_c - 1}$ and wait for its commitment.

\STATE Epilogue: Adjust $\mathbf{O}_{i}$ according to $m_{i}$. Determine $L_{i}$ using both $m_{i}$ and $\ell_{i}$.
      Store $\mathbf{O}_{i}$ and $L_{i}$ in HBM as the $i$-th segments of $\mathbf{O}$ and $L$.
    \end{algorithmic}
  \end{algorithm}

\cref{alg:flash3_wgmma} replaces the consumer path originally presented in \cref{alg:flash3_wgmma_ws_only}, thereby constituting the full \textsc{FlashAttention-3} algorithm tailored for FP16 computation. In summary, we employ WGMMA as a shorthand to denote asynchronous GEMM functionality. In the algorithm’s main loop (see lines \ref{code:mainloop_start} through \ref{code:mainloop_end}), the second invocation of WGMMA in the $j$-th iteration (line \ref{code:second_wgmma}) proceeds in parallel with the softmax computations of the subsequent $(j+1)$-th iteration (line \ref{code:softmax}).

Although the pipelined approach depicted above promises significant theoretical performance benefits, a number of real-world concerns must be addressed:

\paragraph{Compiler reordering} The ideal execution flow illustrated in the pseudocode may not be preserved, as the NVCC compiler frequently rearranges instructions to maximize optimization opportunities. This instruction reordering can interfere with the intended alternation between WGMMA and non-WGMMA operations, potentially undermining performance improvements or introducing unforeseen issues. Nevertheless, examination of the generated SASS code reveals that the desired instruction overlap is achieved as anticipated (see Section~\ref{sec:2-stage-sass}).

\paragraph{Register pressure} Minimizing register spilling is essential for achieving high efficiency. The 2-stage pipeline, however, necessitates allocating additional registers to hold intermediate values and preserve state across pipeline stages. In particular, the algorithm requires keeping a supplementary $\mathbf{S}_{\mathrm{next}}$ in registers, incurring extra usage of $B_r \times B_c \times \text{sizeof}(\text{float})$ per threadblock. This elevated demand for registers may be at odds with increasing block sizes—another standard performance strategy that also intensifies register usage. Effective tuning should therefore be guided by profiling to determine the most beneficial trade-offs.

\paragraph{3-stage pipelining} Building upon the previously discussed 2-stage pipeline, we introduce a 3-stage extension designed to further interleave the second WGMMA computation with the softmax operation. While this method has the potential to push Tensor Core utilization even higher, the inclusion of a third stage exacerbates register requirements, thereby complicating the balance between maximizing tile size and deepening the pipeline. Full algorithmic details and an empirical assessment are provided in ~\cref{sec:3-stage}.

\subsection{Low-precision with FP8}
\label{sec:algofp8}

\textbf{Efficiency: layout transformations.} Executing the forward computation of \textsc{FlashAttention-3} using FP8 precision introduces further difficulties regarding layout compatibility, beyond those present when working with FP16.

Initially, it is important to observe that the tensors $\mathbf{Q}$, $\mathbf{K}$, and $\mathbf{V}$ are usually arranged such that their head dimension is contiguous. However, to adhere to the k-major layout requirement imposed by FP8 WGMMA for the second GEMM, we must ensure that $\mathbf{V}$—specifically, the tiles of $\mathbf{V}$ transferred into SMEM—are contiguous in the sequence length dimension. Given that the TMA load operation cannot alter which dimension is contiguous, we are compelled to either (1) transpose $\mathbf{V}$ in GMEM during a preprocessing phase or (2) perform a tile-wise in-kernel transpose of $\mathbf{V}$ after loading the relevant segments into SMEM. For the preprocessing approach (1), there are two main possibilities: (1a) integrate the transpose operation into the epilogue of an earlier stage such as rotary embedding, thereby fusing the two, or (1b) invoke a dedicated preprocessing transpose kernel\footnote{A high-performance transpose kernel can approach the peak bandwidth of the hardware~\citep{colfax_cutlass_transpose_2024}.} that switches the strides of the head and sequence length axes. Nonetheless, incorporating (1a) into a general-purpose library presents substantial complexity, while (1b) typically results in excessive overhead, particularly under memory-bound conditions characteristic of inference workloads.

Rather, for FP8 \textsc{FlashAttention-3}, we choose approach (2). To perform the in-kernel transpose, we utilize the LDSM (\verb|ldmatrix|) and STSM (\verb|stmatrix|) instructions. These instructions enable a warp of threads to collectively transfer data between SMEM and RMEM, operating on 128-byte segments.\footnote{While the PTX documentation characterizes LDSM/STSM as copying $8 \times 8$ matrices with 16-bit elements \cite[\S 9.7.13.4.15-16]{ptx}, it is possible to combine pairs of 8-bit values to apply LDSM/STSM in the FP8 setting. However, the transpose variants of LDSM/STSM lack the ability to unpack combined 8-bit values, making it necessary to manipulate registers between LDSM and STSM to correctly execute a tile-wise transpose; further specifics are omitted.} These LDSM/STSM instructions offer efficient register usage, which makes them suitable for execution by the producer warpgroup, and they also support layout transposition during memory transfers. Additionally, after the initial iteration, it is feasible to schedule the transposition of the subsequent $\mathbf{V}$ tile during the overlapped execution (in the shadow) of the two WGMMAs that compute using the prior $\mathbf{V}$ and the current $\mathbf{K}$ tile.

Second, we note that, in contrast to FP16, the FP32 accumulator’s memory layout in an FP8 WGMMA does not align with how operand A is stored in registers. \cref{fig:rmem_accum} and \cref{fig:rmem_operand} illustrate segments of these distinct layouts, with entries assigned to thread-specific registers in the presented sequence. Employing byte permute instructions allows us to rearrange the accumulator from the initial WGMMA into the format expected by the subsequent WGMMA, which also aligns with the $\mathbf{V}$ tile arrangement created by the in-kernel transpose. More concretely, referring to \cref{fig:rmem_accum}, we reorder the sequence as
$$\{ \verb|d0 d1 d4 d5 d2 d3 d6 d7| \},$$
repeating this register mapping for every block of 8 bytes. Viewed from the perspective of the $\mathbf{P}$ tile’s logical shape, this operation reorders its columns (e.g., columns $0189$ are repositioned as the first four columns). For WGMMA to produce the correct output tile, the in-kernel transpose can be adjusted to emit a corresponding permutation of the rows in the $\mathbf{V}$ tile. This approach provides extra flexibility by handling the required reordering inside the in-kernel transpose, thereby eliminating the necessity for shuffle instructions to transfer register ownership across threads, as discussed previously in~\citep{colfax_fp8_flashattention_2024}.

\textbf{Accuracy: block quantization and incoherent processing.} In the FP8 (e4m3) representation, the mantissa is encoded with only 3 bits, and the exponent occupies 4 bits. This configuration introduces greater numerical imprecision compared to formats such as FP16 or BF16. Additionally, large-scale models often exhibit outlier values~\citep{dettmers2208llm, sun2024massive} that deviate significantly in magnitude from the bulk of the data, posing challenges for effective quantization. Standard practice involves per-tensor scaling~\citep{micikevicius2022fp8}, whereby a single scaling factor is assigned to each tensor (e.g., individual scalars for $\mathbf{Q}$, $\mathbf{K}$, and $\mathbf{V}$). To mitigate the increased numerical errors of attention computations under FP8, we leverage two strategies:

\begin{enumerate}

\item \textbf{Block quantization}: In this approach, a single scaling factor is retained per block, meaning that the tensors $\mathbf{Q}$, $\mathbf{K}$, and $\mathbf{V}$ are divided into blocks with dimensions $B_r \times d$ or $B_c \times d$. Each block is then quantized independently. This quantization step can be integrated with operations immediately preceding the attention mechanism, such as rotary embedding, without incurring additional computational overhead, since rotary embedding is constrained by memory bandwidth. Given that the \textsc{FlashAttention-3} algorithm inherently processes data in blocks, it is straightforward to rescale each segment of $\mathbf{S}$ to accommodate block quantization without adding computational cost.

\item \textbf{Incoherent processing}: To mitigate the influence of outlier values, both $\mathbf{Q}$ and $\mathbf{K}$ are premultiplied by a random orthogonal matrix $\mathbf{M}$ prior to FP8 quantization. Because $\mathbf{M}$ satisfies $\mathbf{M} \mathbf{M}^\top = I$, it follows that $(\mathbf{Q} \mathbf{M}) (\mathbf{K} \mathbf{M})^\top = \mathbf{Q} \mathbf{K}^\top$, ensuring that the overall attention computation remains unchanged. This procedure distributes the effect of outliers by transforming each element of $\mathbf{Q} \mathbf{M}$ or $\mathbf{K} \mathbf{M}$ into a random linear combination of the original entries, thereby diminishing quantization error. Building on the methods introduced by \citet{chee2024quip} and~\citet{tseng2024quip}, $\mathbf{M}$ is chosen as a composition of random diagonal matrices with entries $\pm 1$ and a Hadamard matrix. This structure allows matrix multiplication in $O(d \log d)$ time (significantly faster than $O(d^2)$) and can be seamlessly merged with the rotary embedding operation without additional computational expense.
\end{enumerate}

Our empirical results, detailed in \cref{sec:numerical_error}, demonstrate that these strategies collectively reduce numerical errors by as much as 2.6$\times$.

\section{Empirical Validation}
\label{sec:exp}

To implement \textsc{FlashAttention-3} and assess its performance and precision, we leverage primitives provided by CUTLASS~\citep{Thakkar_CUTLASS_2023}, including the WGMMA and TMA abstractions.

\begin{itemize}

\item \textbf{Benchmarking attention.} To assess performance, we evaluate the execution time of \textsc{FlashAttention-3} across a variety of sequence lengths and benchmark it against a conventional PyTorch attention implementation, \textsc{FlashAttention-2}, a Triton-based version of \textsc{FlashAttention-2} utilizing H100-specific optimizations, as well as a vendor-provided, cuDNN-accelerated \textsc{FlashAttention-2} optimized for H100 GPUs. Our findings demonstrate that \textsc{FlashAttention-3} achieves speeds up to 2.0$\times$ greater than standard \textsc{FlashAttention-2}, and surpasses Triton-optimized \textsc{FlashAttention-2} by a factor of 1.5$\times$. Furthermore, \textsc{FlashAttention-3} is capable of attaining throughput rates of up to 740 TFLOPs/s, corresponding to 75\% of the peak theoretical TFLOPs/s on H100 hardware.

\item \textbf{Ablation study.} Through systematic ablation, we verify that enhancements such as warp-specialization and GEMM-softmax pipelining play a critical role in the performance gains realized by \textsc{FlashAttention-3}.

\item \textbf{Accuracy of FP8 attention.} We further establish that the adoption of block quantization in conjunction with incoherent processing substantially mitigates the numerical error in FP8 \textsc{FlashAttention-3} by a factor of 2.6$\times$.
\end{itemize}

\subsection{Benchmarking Attention}
\label{subsec:benchmark_attn}

We evaluate the execution time of various attention mechanisms on an H100 80GB SXM5 GPU across multiple configurations, including the presence or absence of a causal mask and head dimensions of either 64 or 128, using FP16 inputs. The findings, presented in~\cref{fig:benchmark_attn_fwd} and~\cref{fig:benchmark_attn_bwd}, indicate that \textsc{FlashAttention-3} achieves a forward pass speedup of approximately 1.5-2.0$\times$ relative to \textsc{FlashAttention-2}, and a 1.5-1.75$\times$ speedup for the backward pass. When benchmarked against a conventional attention implementation, \textsc{FlashAttention-3} attains performance increases of up to 3-16$\times$. Notably, for medium to long sequence lengths (1k and above), \textsc{FlashAttention-3} outperforms even a proprietary and highly optimized vendor library (cuDNN, which is closed source) for H100 GPUs.

\paragraph{Benchmark settings:} We vary the sequence length as 512, 1k, ..., 16k, and set batch size so that the total number of tokens is 16k. We set the hidden dimension to 2048, and head dimension to be either 64, 128, or 256 (i.e., 32 heads, 16 heads, or 8 heads). To calculate the FLOPs of the forward pass, we use:

\begin{equation*}
  4 \cdot \text{seqlen}^2 \cdot \text{head dimension} \cdot \text{number of heads}.
\end{equation*}

When applying causal masking, we halve this value, since causal masking ensures that only about fifty percent of the elements need to be computed. To estimate the FLOPs for the backward pass, we multiply the FLOPs from the forward pass by 2.5; this scaling factor reflects that the backward computation involves five matrix multiplications compared to two in the forward computation (owing to recomputation during backpropagation).

In addition, we evaluate the forward pass runtime for FP8 under comparable experimental conditions. The outcomes corresponding to a headdim of 256 are illustrated in~\cref{fig:benchmark_attn_fp8}, while comprehensive findings are provided in \cref{sec:benchmark_attn_fp8_full}.

\subsection{Ablation Study: 2-Stage Pipelining Experiments}
\label{sec:2-stage-pipelining-experiments}

We examine the effects of removing both the 2-stage WGMMA-softmax pipeline and warp-specialization in non-causal FP16 \textsc{FlashAttention-3} under fixed parameters $\{\text{batch}, \text{seqlen}, \text{nheads}, \text{hdim}\} = \{ 4, 8448, 16, 128\}$. Findings presented in~\cref{table:ablation_pipelining} demonstrate that our enhancements—namely, introducing asynchrony via warp-specialization and enabling the concurrent execution of GEMM and softmax—yield a marked performance boost, increasing throughput from 570 to 661 TFLOPs.

\subsection{Numerical Error Validation}
\label{sec:numerical_error}

Given recent concerns regarding the numerical error~\citep{golden2024flash} present in \textsc{FlashAttention}, we evaluate \textsc{FlashAttention-2}, \textsc{FlashAttention-3}, and a conventional attention implementation, using an FP64 reference as a baseline. To replicate the outlier features and activations typically observed in LLMs~\citep{dettmers2208llm, sun2024massive}, we sample the elements of $\mathbf{Q}, \mathbf{K}, \mathbf{V}$ according to the following distribution:

\begin{equation*}
  \mathcal{N}(0, 1) + \mathcal{N}(0, 100) \cdot \mathrm{Bernoulli}(0.001).
\end{equation*}

Specifically, each matrix entry follows a normal distribution centered at zero with unit standard deviation, except for 0.1\% of the entries, which are modified by adding a separate, independent normally distributed value with standard deviation 10. The root mean squared error (RMSE) results are presented in~\cref{table:numerical_error}. When operating in FP16, both \textsc{FlashAttention-2} and \textsc{FlashAttention-3} achieve a 1.7$\times$ reduction in RMSE relative to the standard approach, due to maintaining intermediate computations (specifically, the softmax output) in FP32 precision. The FP8 baseline employs per-tensor scaling, utilizes FP32 for matmul accumulation, and keeps softmax intermediates in FP16. Owing to block quantization and its incoherent computation scheme, \textsc{FlashAttention-3} operating in FP8 demonstrates a 2.6$\times$ improvement in accuracy over the baseline.

\section{Dicussion, Limitations, Conclusion}
\label{sec:discussion}

Through the development of \textsc{FlashAttention-3}, we have illustrated how advancements in programming methods and the incorporation of hardware capabilities—specifically asynchrony and low-precision computation—can significantly enhance both the efficiency and the reliability of attention mechanisms. Our approach accelerates attention by a factor of 1.5 to 2.0$\times$ compared to \textsc{FlashAttention-2}, and decreases FP8 numerical inaccuracies by 2.6$\times$ relative to conventional per-tensor quantization techniques. However, several challenges remain. Future work will focus on optimizing performance for LLM inference, adopting a persistent kernel architecture within the FP8 kernel,\footnote{In our evaluations, the FP16 implementation of \textsc{FlashAttention-3} utilizes both a persistent kernel and load balancing strategy, unlike the FP8 implementation. This discrepancy partially accounts for the reduced performance of FP8 \textsc{FlashAttention-3} on shorter sequence lengths and under causal masking conditions, particularly when compared to the FP8 cuDNN kernels.} and investigating how low-precision attention affects large-scale model training. While this study primarily targets Hopper GPUs, we anticipate that the underlying techniques can generalize to additional hardware accelerators. Ultimately, we aspire that a more efficient and precise attention module will facilitate the development of novel long-context applications.

\end{document}